\section{برش دستی چرم نازک و ضخیم}

انتخاب ابزار به ضخامت، سفتی و شکل مسیر بستگی دارد. چرم نازک و نرم ممکن است زیر تیغ کشیده شود. در چرم ضخیم و سفت، فشار زیاد در یک عبور می‌تواند تیغ را از مسیر یا از راستای عمود منحرف کند. هر روش را ابتدا روی دورریز همان پوست، با همان جهت برش و زیرانداز، آزمایش کنید.

\begin{table}[!htbp]
  \centering
  \caption{انتخاب ابزار برای شکل‌های مختلف برش دستی}
  \label{tab:tools}
  \begin{tabular}{@{}p{30mm}p{52mm}p{64mm}@{}}
    \toprule
    مسیر & ابزار قابل‌آزمون & کنترل اصلی \\
    \midrule
    خط مستقیم & خط‌کش فلزی پایدار و تیغ مناسب. کاتر گرد برای بعضی نمونه‌های نرم & خط‌کش روی بخش مفید قطعه باشد. موقعیت لبه را در رویه و پشت کنترل کنید. \\
    منحنی باز & کاتر گرد برای منحنی ملایم یا تیغ ظریف با حرکت‌های کوتاه & پوست را نکشید و ابزار را در امتداد قوس حرکت دهید. \\
    گوشۀ داخلی و شکاف & تیغ نوک‌تیز و، در صورت وجود گردیِ طراحی‌شده، پانچ با قطر تعیین‌شده & دو خط را در نقطۀ پایان به هم برسانید. عبور از گوشه درز را ضعیف می‌کند. \\
    سوراخ یا شیار & پانچ گرد یا پانچ شیار با زیرانداز سازگار & از مرکز، قطر و فاصلۀ سوراخ تا لبه نمونه بگیرید. قطر را با یراق واقعی مقایسه کنید. \\
    نوار بلند & خط‌کش یا بندبُر روی یک لبۀ مرجعِ صاف & پهنا را در ابتدا، میانه و انتها کنترل کنید. \\
    \bottomrule
  \end{tabular}
\end{table}

جدول~\ref{tab:tools} نقطۀ شروع انتخاب است. برای کاتر گرد، از زیرانداز و خط‌کش سازگار استفاده کنید و تیغ را در زمان توقف بپوشانید. خط برش را دور از دست نگه‌دارنده و در راستایی بگیرید که ابزار به‌سمت بدن حرکت نکند \cite{olfa}.

\subsection{چرم نازک و نرم}

چرم را روی زیرانداز صاف کنید و در حالت آزاد نگه دارید. وزنه‌های نرم یا راهنمای ضدلغزش باعث می‌شوند برای ثابت‌کردن قطعه به چرم کشش ندهید. اگر تیغ در دورریز چین می‌اندازد، ابتدا تیزی، نوع تیغ و تکیه‌گاه را بررسی کنید. کاتر گرد روی مسیر باز می‌تواند کشیده‌شدنِ سطح را کمتر کند. در گوشۀ بسته، تیغ ظریف معمولاً امکان توقف دقیق‌تری می‌دهد.

\subsection{چرم ضخیم و سفت}

روی دورریز بررسی کنید که لبه پس از برش از سمت رو و پشت در یک راستا باشد. برای مسیر مستقیم، می‌توان عبور اول را سبک و دقیق گرفت و عمق را در چند عبور کنترل‌شده افزایش داد. اگر تیغ شروع به کج‌شدن یا گیرکردن کرد، برش را متوقف کنید و تیزی ابزار، وضعیت زیرانداز و زاویۀ تیغ را بررسی کنید. برای قطعات جفت، پهنای قطعۀ دوم را با الگو بسنجید تا خطای قطعۀ اول تکرار نشود.

\subsection{ترتیب اجرای برش}

تا زمانی که قطعۀ خام بزرگ و پایدار است، سوراخ‌ها و پنجره‌های داخلیِ لازم را اجرا کنید. اگر محل سوراخ به مونتاژ یا یراق وابسته است، مرکز را علامت بزنید و سوراخ‌کاری را به پس از کنترل تطبیق موکول کنید. گوشه‌ها را با دو حرکت کنترل‌شده به نقطۀ مشترک برسانید. در پایان، محیط قطعه را ببرید و اندازه‌ها، گونیا، کیفیت منحنی و رنگ یا آسیب لبه را با شابلون مقایسه کنید. برای برش چندلایه، اول بررسی کنید که الیاف و لایه‌های زیرین در آزمون سر نمی‌خورند و برش تا انتها کامل است.

\section{لویس‌کاری و کنترل ضخامت}

لویس‌کاری (\lr{Skiving}) کم‌کردن ضخامت در بخشی از قطعه است. نازک‌کردنِ تدریجیِ لبه برای هم‌پوشانی، و نازک‌کردن روی خط تا برای کم‌کردن برآمدگی، دو کاربرد رایج‌اند. لایه‌بری (\lr{Splitting}) می‌تواند ضخامت ناحیۀ وسیع‌تری از قطعه را عوض کند. نقش این دو عملیات را پیش از برش نهایی از هم مشخص کنید \cite{sewpro}.

\begin{figure}[!htbp]
  \centering
  \SkiveDiagram
  \caption{مقطع شماتیک لویس لبه. علامت \lr{A} ضخامت اصلی، \lr{B} پهنای گذار و \lr{C} ضخامت باقی‌مانده در لبه را نشان می‌دهد.}
  \label{fig:skive}
\end{figure}

در شکل~\ref{fig:skive}، رویه صاف مانده و نازک‌کاری از پشت انجام شده است. روی دورریز ابتدا پهنای لویس، ضخامت آغازین و ضخامت باقی‌مانده را ثبت کنید. کولیس ممکن است چرم نرم را فشرده کند. برای مقایسۀ کارگاهی ابزار و فشار را یکسان نگه دارید. برای گزارش آزمایشگاهی روش استاندارد تعیین ضخامت چرم را مبنا بگیرید \cite{iso-thickness}.

\subsection{لویس دستی و ماشینی}

در روش دستی، پشت چرم را روی سطح تخت و مناسب ابزار بگذارید و با تیغ تیز، برداشت کم و پی‌درپی انجام دهید. اگر پهنای گذار ثابت نمی‌ماند، خط مرز را دوباره کنترل کنید. برای لبه‌های باریک یا محل دوخت، باقی‌ماندن ضخامت کافی اهمیت دارد. قبل از تکرار کار روی قطعات اصلی، یک نمونۀ چندلایه را بچسبانید، خم کنید و بدوزید.

در دستگاه تیغ کاسه‌ای (\lr{Bell-knife skiver})، غلتک چرم را به ناحیۀ برش می‌رساند. در نمونه‌ای مانند \lr{SewPro 801}، راهنمای لبه پهنای لویس و کفشک ضخامت باقی‌مانده را کنترل می‌کند. نوع غلتک و تنظیم قابل‌استفاده به مدل و جنس قطعه وابسته است. از تکه‌ای با ضخامت و جهت واقعی شروع کنید و هر بار فقط یک متغیر را تغییر دهید. حفاظت تیغ و روش تنظیم یا سرویس را از دفترچۀ مدل دستگاه بگیرید \cite{sewpro,guard}.

کَنده‌شدن لبه، نازک‌شدن خط دوخت یا موج‌دارشدن پشت، نشانه‌ای است که نمونه باید دوباره بررسی شود. محل توقف، فشار تغذیه، تیزی و سازگاری دستگاه با قطعه را در چارچوب دستور سازنده بررسی کنید. قطعه‌ای را که رویه‌اش بریده یا محل دوختش ضعیف شده برای اتصال باربر استفاده نکنید.

\section{پرس و قالب تیغ‌دار}

قالب تیغ‌دار (\lr{Cutting die}) با اعمال نیرو از پرس، مرز قطعه را می‌بُرد. شابلون معمولی خط را منتقل می‌کند. پرسِ سوراخ‌کاری، نصب یراق یا نقش‌زنی ممکن است ابزار و شرایط کاری دیگری داشته باشد. پیش از استفاده، سازگاری قالب، صفحه و دستگاه را با کاربرد مدنظر از مشخصات سازنده بررسی کنید \cite{tippmann,weaver}.

\begin{figure}[!htbp]
  \centering
  \PressStackDiagram
  \caption{چیدمان شماتیک پرس برش. علامت‌های \lr{A} تا \lr{D} به‌ترتیب صفحۀ فشار، قالب تیغ‌دار، چرم و صفحۀ برش‌اند. جهت و آرایش واقعی تابع دفترچۀ دستگاه است.}
  \label{fig:press}
\end{figure}

شکل~\ref{fig:press} فقط ترتیب اجزا را در یک چیدمان نمونه نشان می‌دهد. پیش از برش، جاگرفتن کامل قالب در سطح کاری، ارتفاع قابل‌پذیرش، ضخامت و جنس چرم، و اندازۀ صفحۀ محافظ را بررسی کنید. پرس دستی برای کار مقطعی و قالب‌های سازگار مناسب است. برای تکرار بالاتر، ثبات چرخه، عرض مفید و نیازهای تأمین هوا یا سرویس را در کنار نوع دستگاه بررسی کنید. نیروی اسمی به‌تنهایی ظرفیت برش قالبِ مشخص را تعیین نمی‌کند \cite{tippmann,weaver}.

خطر اصلی در ناحیۀ برخورد قالب و صفحه است. هنگام آغاز چرخه، دست را بیرون از ناحیۀ برش نگه دارید و حفاظ یا کنترل تعریف‌شده برای دستگاه را به‌کار ببرید. قبل از جابه‌جایی قطعه، بازگشت کامل صفحه را بررسی کنید. هرگونه سرویس یا رفع گیر باید با دستور توقف و جداسازی انرژیِ همان دستگاه انجام شود \cite{tippmann,guard}.

\subsection{برش آزمایشی و تشخیص خطا}

یک تکه از همان سری چرم را آزمایش کنید. جدایش کامل دور قطعه، برش گوشه‌ها و وضعیت سطح را در هر دو طرف ببینید. اگر بخشی از برش کامل نیست، اول سالم‌بودن تمام تیغه، یکنواختی صفحۀ برش، قرارگیری قالب، ضخامت واقعی و پوشش کامل صفحه را بررسی کنید. در قطعه‌ای که لبه‌اش له شده یا رویه‌اش اثر فشار دارد، سلامت صفحه، تیزی قالب و تناسب روش را بررسی کنید. تغییر فشار یا تنظیم محدودۀ حرکت فقط در حدود و با روش تعیین‌شده در دفترچۀ سازنده انجام می‌شود \cite{tippmann}.
